% 妊神星（综述）
% license CCBYSA3
% type Wiki

本文根据 CC-BY-SA 协议转载翻译自维基百科\href{https://en.wikipedia.org/wiki/Haumea}{相关文章}。

Haumea（小行星编号：136108 Haumea）是一颗位于海王星轨道之外的矮行星。$^{[25]}$
它由加州理工学院（Caltech）的 Mike Brown 领导的团队于 2004 年在帕洛马天文台发现；随后，西班牙内华达山脉天文台 José Luis Ortiz Moreno 领导的团队在 2005 年正式宣布了这一发现，他们在 2003 年拍摄的先前图像中识别出了该天体。从那次宣布起，它获得了临时编号 2003 EL61。

2008 年 9 月 17 日，它被正式命名为 Haumea，源自夏威夷生育与分娩女神。国际天文学联合会（IAU）基于其可能具备矮行星地位的预期而赋予这一名称。根据标称估计，它是仅次于厄里斯和冥王星的第三大已知海王星外天体，其大小约等同于天王星的卫星泰坦妮娅。Haumea 的最早“预发现”（precovery）影像可以追溯至 1955 年 3 月 22 日。$^{[9]}$

Haumea 的质量约为冥王星的三分之一、地球的 1/1400。虽然尚未直接观测到其形状，但根据其光变曲线的计算结果，与 雅可比椭球体（Jacobi ellipsoid）一致（这是一种在满足矮行星条件时会呈现的形状），其长轴大约是短轴的两倍。

2017 年 10 月，天文学家宣布在 Haumea 周围发现了一套环系统——这是首个在海王星外天体以及矮行星上发现的环系统。

Haumea 的引力直到最近仍被认为足以使其达到流体静力平衡，但目前这一点仍不确定。Haumea 修长的形状、快速自转、环结构以及高反照率（源自其晶质水冰表面），被认为都是一次巨大碰撞的结果。这场碰撞使 Haumea 成为一个碰撞族群（Haumea 家族）的最大成员，该族群包括多个大型海王星外天体，以及 Haumea 的两个已知卫星：Hiʻiaka 与 Namaka。
\subsection{历史}
\subsubsection{发现}
由加州理工学院（Caltech）的 Mike Brown、耶鲁大学的 David Rabinowitz，以及位于夏威夷双子天文台的 Chad Trujillo 组成的团队，于 2004 年 12 月 28 日在他们拍摄的2004 年 5 月 6 日的图像中发现了 Haumea。2005 年 7 月 20 日，他们在网上发布了一篇摘要，准备在 2005 年 9 月的一次会议上正式宣布这一发现。$^{[26]}$

在此期间，西班牙内华达山脉天文台安达卢西亚天体物理研究所（Instituto de Astrofísica de Andalucía）的José Luis Ortiz Moreno及其团队，在2003 年 3 月 7–10 日拍摄的图像中找到了 Haumea。2005 年 7 月 27 日晚，Ortiz 向小行星中心（MPC）发送了关于他们发现的邮件。$^{[27]}$

Brown 最初承认发现权应归 Ortiz 团队，$^{[28]}$ 但在得知西班牙团队在官方发现宣布的前一天访问过 Brown 的观测日志后，开始怀疑对方存在欺诈行为，而 Ortiz 团队在宣布中并未像惯例那样披露这一事实。这些日志包含足够的数据，使 Ortiz 团队能够在其 2003 年的图像中提前定位（precover）Haumea，并且在 Ortiz 在7 月 29 日安排望远镜时间以获取确认图像、并向 MPC 再次提交发现报告之前，又再次访问了这些日志。Ortiz 后来承认访问过 Caltech 的观测日志，但否认任何不当行为，声称自己只是核实他们是否发现了一个新的天体。$^{[29]}$

根据 IAU（国际天文学联合会）规定，小行星的发现权归属于最先向 MPC（小行星中心）提交足够轨道数据以确定其轨道的团队，并且获得发现权的团队有权优先命名该天体。
然而，IAU 在2008 年 9 月 17 日发布的关于 Haumea 命名的官方公告（该命名由专门负责预期为矮行星的天体设立的联合委员会决定）中，并未提及发现者。公告中将发现地点列为西班牙团队的内华达山脉天文台，$^{[30][31]}$ 但最终采用的名称“Haumea”却是 Caltech 团队的提议。

Ortiz 团队曾提议命名为“Ataecina”，取自古伊比利亚春之女神；$^{[27]}$ 但由于该女神属于冥界神（chthonic deity），这一名字更适合用于命名“冥卫星类（plutino）”天体，而 Haumea 并非此类天体，因此未被采用。
\subsubsection{名称与符号}
在被正式命名之前，加州理工学院（Caltech）发现团队内部给 Haumea 起了绰号“Santa（圣诞老人）”，因为他们是在2004年12月28日，也就是圣诞节刚过后发现该天体的。$^{[32]}$

西班牙团队于2005年7月首次向小行星中心（MPC）提交了发现申请。2005 年 7 月 29 日**，Haumea 根据西班牙团队发现图像的日期，获得了临时编号2003 EL61。2006 年 9 月 7 日**，它被编号并正式列入小行星目录，成为(136108) 2003 EL61。

根据当时国际天文学联合会（IAU）规定：经典柯伊伯带天体应以与创造相关的神话人物命名，$^{[33]}$ 因此 2006 年 9 月，Caltech 团队向 IAU 正式提交了来自夏威夷神话的名称建议，用于命名 (136108) 2003 EL61 及其卫星，以“向卫星发现地致敬”。$^{[34]}$ 这些名称由 Caltech 团队成员 David Rabinowitz 提出。$^{[25]}$

Haumea 是夏威夷岛（Hawaiʻi）的守护女神，双子座天文台（Gemini）与凯克天文台（W. M. Keck Observatory）就坐落在冒纳凯亚山（Mauna Kea）。此外，她也被视作大地女神 Papa，是天空之神 Wākea（象征太空）的妻子。$^{[35]}$ 在当时，这一命名被认为非常贴切，因为当时 Haumea 被认为几乎完全由岩石构成，而不像其他柯伊伯带天体那样拥有“岩石核心 + 厚冰幔”的结构。$^{[36][37]}$

此外，Haumea 是 生育与分娩女神，据说许多子女从她身体的不同部位诞生；$^{[35]}$ 这与理论推测相对应——Haumea 主体在远古一次碰撞中碎裂，形成了大量冰质天体（即所谓 Haumea 碰撞家族）。$^{[37]}$ Haumea 已知的两颗卫星也被认为是在这一碎裂事件中形成，因此以其两个女儿 Hiʻiaka 与 Nāmaka 命名。$^{[36]}$

西班牙 Ortiz 团队提出的命名 Ataecina 不符合 IAU 命名规定，因为冥界神（chthonic deity）的名字仅用于稳定共振的柯伊伯带天体，例如与海王星呈 3:2 共振的冥卫星类（plutinos），而 Haumea 处于不稳定的 7:12 间歇共振状态，因此严格意义上并非共振天体。此命名规则在 2019 年末 被 IAU 明确：冥界神名称今后专属用于冥卫星类天体。

Haumea 的行星符号为 ⟨🝻⟩，现收录于 Unicode 编码 U+1F77B。$^{[38]}$行星符号在现代天文学中已较少使用，🝻 主要由占星师使用，$^{[39]}$ 但 NASA 也曾使用过该符号。$^{[40]}$
该符号由马萨诸塞州软件工程师 Denis Moskowitz 设计，其造型结合并简化了夏威夷岩画中代表“女性”和“分娩”的符号。$^{[41]}$
\subsection{轨道}
Haumea 的轨道周期为 284 个地球年，近日点距离约 35 天文单位（AU），轨道倾角为 28°。$^{[9]}$ 它在 1992 年初经过远日点，目前距离太阳已经超过 50 AU。$^{[23]}$ Haumea 将在 2133 年抵达近日点。$^{[10]}$

Haumea 的轨道偏心率略高于其碰撞家族其他成员。人们认为，这主要是因为 Haumea 与海王星之间存在一个较弱的 7:12 轨道共振，在大约十亿年的时间尺度上逐渐改变了其原始轨道，$^{[37][42]}$ 这种演化过程受 Kozai 效应（古在效应）影响，该效应允许轨道倾角与偏心率之间发生交换。$^{[37][43][44]}$

Haumea 的视星等为 17.3，$^{[23]}$ 是柯伊伯带中继冥王星与鸟神星（Makemake）之后的第三亮天体，使用大型业余望远镜即可轻松观测。$^{[45]}$

然而，由于行星与大多数小型太阳系天体都形成于原始太阳星云盘，因此它们具有共同的轨道平面投影，即黄道面（ecliptic）。$^{[46]}$ 早期对远距天体的搜索，大多集中在天空上接近黄道面的区域。当这一地区逐渐被充分勘测后，后续的天空巡天开始寻找轨道倾角更高、以及更遥远、运动更缓慢的天体。$^{[47][48]}$ 正是在这些扩展后的巡天中，才覆盖到 Haumea 所在的天空位置，因此得以发现该天体，因为它具有较高的轨道倾角，并且目前位于远离黄道的位置。
\subsubsection{可能的与海王星轨道共振}
哈乌梅亚被认为处于与海王星的间歇性 7:12 轨道共振$^{[37]}$。其升交点 $\Omega$ 的进动周期约为 460 万年，每一个进动周期共振会中断两次（即每约 230 万年一次），在约十万年后又重新恢复$^{[5]}$。由于这种共振并不稳定，Marc Buie 将其归类为非共振状态$^{[49]}$。
\subsection{自转}
哈乌梅亚在约 3.9 小时的周期内呈现出剧烈的亮度变化，这只能用其自转周期来解释$^{[50]}$。这比太阳系中任何其他已知达到流体静力平衡（equilibrium body）的天体都更快，事实上也比任何直径超过 100 公里的已知天体自转得更快$^{[45]}$。大多数达到平衡形状的自转天体会被压扁成扁球体（oblate spheroid），但哈乌梅亚自转速度极快，因此被拉伸成一个三轴椭球体。如果哈乌梅亚再快一些，它将变形为“哑铃形”，最终可能裂成两个天体$^{[25]}$。这种快速自转被认为是形成其卫星及碰撞族群的大撞击事件所造成的$^{[37]}$。

目前从地球观察，哈乌梅亚的赤道平面几乎是边缘视角（edge-on），并且相对于其环轨道和最外侧卫星希伊阿卡（Hiʻiaka）的轨道平面略有偏移。虽然 Ragozzine 和 Brown 在 2009 年最初假设二者是共面，但他们关于哈乌梅亚卫星碰撞形成模型的研究一直表明，哈乌梅亚赤道平面至少与希伊阿卡的轨道平面对齐，差异约为 $1^\circ$.$^{[15]}$

这一观点在 2017 年通过一次恒星掩食观测得到支持，该观测揭示了哈乌梅亚存在一个环，其平面大致与希伊阿卡的轨道平面及哈乌梅亚赤道相重合$^{[12]}$。Kondratyev 与 Kornoukhov 在 2018 年对掩食数据进行的数学分析进一步限制了哈乌梅亚赤道相对于其环平面和希伊阿卡轨道平面的相对倾角，分别为 $3.2^\circ \pm 1.4^\circ$ 与 $2.0^\circ \pm 1.0^\circ$，均相对于哈乌梅亚赤道$^{[17]}$。
\subsection{物理特征}
\subsubsection{大小、形状与成分}
太阳系天体的大小可以根据其光学星等、距离以及反照率来推断。天体对地球观测者而言显得明亮，要么是因为它体积大，要么是因为它高度反射。如果能够确定其反照率（albedo），就可以对它的大小做一个粗略估计。对于大多数遥远天体来说，其反照率是未知的，但 Haumea 足够巨大且明亮，可以测量其热辐射，从而获得一个关于其反照率以及大小的近似值$^{[51]}$。然而，由于 Haumea 的快速自转，其尺寸计算变得复杂。可变形天体的自转物理预测，在短短约一百天内$^{[45]}$，像 Haumea 这样快速自转的天体会被拉伸为一个平衡形态的三轴椭球。一般认为，Haumea 光度曲线的大部分波动，并非由局部反照率差异引起，而是由于从地球观察时视线在侧面视图与端面视图之间交替变化$^{[45]}$。

Haumea 的自转和光变幅度被认为对其成分施加了强有力的约束。如果 Haumea 处于流体静力平衡，并且像冥王星一样具有低密度，即一个小型岩质核心外覆盖厚厚的冰层，那么其快速自转将导致其形状被拉长到一个比光变所允许的程度更高的数值。此类分析将其密度限制在 $2.6-3.3 g/cm^3$ 的范围内$^{[52][45]}$。相比之下，呈岩质的月球密度为 $3.3 g/cm^3$，而典型柯伊伯带冰质天体冥王星的密度为 $1.86 g/cm^3$。Haumea 可能的高密度覆盖了许多太阳系岩质天体所具有的硅酸盐矿物（如橄榄石和辉石）的密度值。这也暗示 Haumea 的主体结构是岩石，并覆盖有一层相对较薄的冰层。典型柯伊伯带天体所具备的厚冰外壳，可能在形成 Haumea 撞击族的巨大碰撞中被剥离$^{[37]}$。

由于 Haumea 具有卫星，可以利用开普勒第三定律根据这些卫星的轨道计算 Haumea 系统的质量。结果为 $4.2 \times 10^{21}$ kg，为冥王星系统质量的 28\%，约为月球质量的 6\%。几乎所有这些质量都集中在 Haumea 本身$^{[15][53]}$。

针对 Haumea 大小的多个椭球模型已被提出。Haumea 被发现之后，通过地基望远镜对其光度曲线进行光学波段测量的第一个模型，给出了 1960–2500 km 的总长度，以及大于 0.6 的可见光反照率（$p_v$）$^{[45]}$。最可能的形状是一个近似 2000 × 1500 × 1000 km 的三轴椭球，反照率约为 0.71$^{[45]}$。斯皮策空间望远镜（Spitzer）在 70 μm 红外波段光度测量中给出了 $1150^{+250}*{-100}$ km 的直径，以及 $0.84^{+0.1}*{-0.2}$ 的反照率$^{[51]}$。后续的光变曲线分析提出了约 1450 km 的等效圆直径$^{[54]}$。2010 年，通过赫歇尔空间望远镜与早期斯皮策望远镜数据的联合分析，提出了 Haumea 约 1300 km 的新的等效直径$^{[55]}$。这些独立的大小估计在大约 1400 km 的几何平均直径附近相互重叠。2013 年，赫歇尔空间望远镜测得 Haumea 的等效圆直径约为 $1240^{+69}_{-58}$ km$^{[56]}$。

然而，2017 年 1 月一次恒星掩食观测对上述所有结论提出了质疑。虽然 Haumea 的形状如此前推测般呈拉长状态，但测量结果显示其尺寸明显更大——根据掩食观测数据，Haumea 沿其最长轴的直径大约相当于冥王星，而两极方向则约为前者的一半$^{[12]}$。由 Haumea 实测形状计算得到的密度约为 1.8 g/cm$^3$，更符合其它大型 TNO（跨海王星天体）的密度范围。这一结果所对应的形状似乎与均质、处于流体静力平衡状态的天体不一致$^{[12]}$，尽管如此，Haumea 仍然是目前发现的最大跨海王星天体之一$^{[51]}$，大小仅次于 Eris 和 Pluto，与 Makemake 相近，也可能与 Gonggong 类似，并且大于 Sedna、Quaoar 和 Orcus。

2019 年的一项研究尝试通过将 Haumea 设定为分化天体的数值模型来解决有关其形状和密度的相互矛盾测量。研究发现，≈ 2,100 × 1,680 × 1,074 km（长轴按 25 km 步长建模）的尺寸最符合 2017 年掩食观测到的 Haumea 形状，同时也与处于流体静力平衡状态的表层与核心不等轴椭球形状相一致$^{[11]}$。新的形状解暗示 Haumea 的核心尺寸约为 1,626 × 1,446 × 940 km，密度较高，约 ≈ 2.68 g/cm$^3$，指示其主要由水合硅酸盐（如高岭石）构成。该核心被冰质地幔包围，地幔厚度在两极约 70 km，而在最长轴方向约为 170 km，占 Haumea 总质量的最多 17\%。Haumea 的平均密度估计约为 $\approx2.018 g/cm^3$，反照率约为 $\approx0.66$ $^{[11]}$。
\subsubsection{表面}
2005 年，双子座（Gemini）与 Keck 望远镜获得的 Haumea 光谱显示其表面呈现出强烈的晶质水冰特征，与冥王星卫星 Charon 的表面类似$^{[20]}$。这颇为反常，因为晶质冰形成于高于 110 K 的温度，而 Haumea 表面温度低于 50 K，在此温度下应形成无定形冰$^{[20]}$。此外，在跨海王星天体上，晶质冰结构会在太阳高能粒子与宇宙射线持续轰击下变得不稳定$^{[20]}$。在这种 bombardment（轰击）下，晶质冰退化为无定形冰的时间尺度约为一千万年$^{[57]}$，然而跨海王星天体在其现有的低温环境中已存在了数十亿年$^{[42]}$。

辐射损伤还会使跨海王星天体表面在存在有机冰以及类似 tholin 的物质时变红、变暗，这在冥王星表面就很常见。因此，光谱和颜色表明 Haumea 及其族群天体经历了近期的再表面化过程（resurfacing），从而产生了新鲜的冰。然而，目前尚未提出合理的再表面化机制$^{[22]}$。

Haumea 明亮如雪，反照率在 0.6–0.8 范围内，与晶质冰一致$^{[45]}$。其它大型 TNO，如 Eris，其反照率可能与此一样高甚至更高$^{[58]}$。对表面光谱的最佳拟合模型显示，Haumea 表面约有 66\%–80\% 似乎由纯晶质水冰组成，导致其高反照率的一个可能贡献物是氰化氢或页硅酸盐黏土$^{[20]}$。无机氰化物盐，如铜钾氰化物，也可能存在$^{[20]}$。

然而，对可见光及近红外光谱的进一步研究显示，其表面均一，覆盖着约 1:1 比例的无定形与晶质冰混合物，同时含不超过 8\% 的有机物。无氨水合物的存在排除了低温火山活动（cryovolcanism），这些观测结果也证实，Haumea 的撞击事件必须发生在一亿年以上前，与动力学研究一致$^{[59]}$。在 Haumea 光谱中未检测到可测量的甲烷含量，这与一次会去除挥发物的温暖撞击历史相一致$^{[20]}$，这与 Makemake 的情况形成对比$^{[60]}$。

除了 Haumea 光度曲线因其自身形状而产生的巨大波动（对所有颜色均等影响）外，在可见光与近红外波段还检测到较小、独立的颜色变化，这表明其表面存在颜色与反照率不同的区域$^{[61][62]}$。更具体地说，在 2009 年 9 月，在 Haumea 明亮的白色表面上观测到一个巨大的暗红区域，可能是撞击产物，表明该区域富含矿物与有机（含碳）化合物，或可能含有更高比例的晶质冰$^{[50][63]}$。因此，Haumea 的表面可能呈现出类似冥王星的杂色外观，尽管不如后者那样极端。
\subsection{环}
恒星掩食观测到的环结构一项于 2017 年 1 月 21 日观测到的恒星掩食，并在 2017 年 10 月发表于 《Nature》 的研究中进行了描述，显示海卫密（Haumea）存在一条行星环。这是首次在跨海王星天体（TNO）中发现环状系统$^{[12][64]}$。该环半径约为 2,287 km，宽度约为 70 km，不透明度约为 0.5。该环位置完全处于海卫密的洛希极限之内；如果海卫密是球形，其洛希极限半径大约为 4,400 km（由于海卫密并非球形，这一极限会被推得更远）$^{[12]}$。

该环平面相对于海卫密赤道平面的倾角为 $3.2^\circ \pm 1.4^\circ$，并与其较大、外层卫星希伊阿卡（Hiʻiaka）的轨道平面大致重合$^{[12][65]}$。该环的位置也非常接近海卫密自转的 1:3 轨道–自转共振（其中心距海卫密中心的半径为 $2,285 \pm 8 \text{ km}$）。据估计，该环约贡献了海卫密总光度的 5\%$^{[12]}$。

在 2019 年发表的一项关于环粒子动力学的研究中，Othon Cabo Winter 及其同事表明，与海卫密自转的 1:3 共振在动力学上是不稳定的，但在相空间中存在一处稳定区域，与海卫密环所处位置一致。这表明环中微粒源于接近但未进入该共振区的圆形周期轨道$^{[66]}$。
\subsection{卫星}
两颗已发现的卫星已发现有两颗小卫星环绕海卫密运行，即 (136108) Haumea I，命名为希伊阿卡（Hiʻiaka），以及 (136108) Haumea II，命名为 纳玛卡（Namaka）$^{[30]}$。这两颗卫星由达林·拉戈津（Darin Ragozzine）和迈克尔·布朗（Michael Brown）于 2005 年在使用 W. M. 凯克天文台（W. M. Keck Observatory）观测海卫密的过程中发现。

希伊阿卡（Hiʻiaka）最初被加州理工团队昵称为“鲁道夫（Rudolph）”$^{[67]}$，于 2005 年 1 月 26 日被发现$^{[53]}$。它是外层卫星，直径约 310 km，是两颗卫星中较大且更明亮的一个，每约 49 天以近圆轨道绕海卫密一周$^{[68]}$。其红外光谱在 1.5 微米和 2 微米处存在强烈吸收特征，与几乎完整覆盖其表面的高纯度晶体水冰一致$^{[69]}$。这一不寻常的光谱特征，加上海卫密本体上观察到的类似吸收线，使布朗及其同事得出结论：“俘获形成模型”不太可能；海卫密卫星系统应源自海卫密自身的碎裂物$^{[42]}$。

纳玛卡（Namaka）是海卫密的内层、较小卫星，于 2005 年 6 月 30 日被发现$^{[70]}$，昵称为“闪电（Blitzen）”。其质量约为希伊阿卡的 1/10，以高度椭圆、非开普勒轨道，每 18 天绕海卫密一周，并截至 2008 年，其轨道与外层卫星存在约 13° 的倾角，从而受到后者的显著摄动$^{[71]}$。两颗卫星轨道的较大偏心率及相互倾角是出乎预料的，因为潮汐效应本应对其进行阻尼。一项较新的解释认为：海卫密卫星系统可能在近期经历过 3:1 轨道共振，从而产生当前被“激发”的轨道状态$^{[72]}$。

自约 2008 年至 2011 年$^{[73]}$，从地球观测，海卫密卫星系统的轨道几乎呈边缘朝向，其中纳玛卡会周期性地掩食海卫密$^{[74]}$。观测此类凌掩现象将可获得海卫密及其卫星的精确尺寸与形状信息$^{[75]}$，类似于 20 世纪 80 年代末对冥王星及其卫星卡戎（Charon）的观测$^{[76]}$。不过，由于这些掩食过程中系统亮度的微小变化很难检测，至少需要中等口径的专业望远镜$^{[75][77]}$。希伊阿卡上一次掩食海卫密是在 1999 年，比发现时间早几年，并将在约 130 年后才再次发生$^{[78]}$。

然而，在这一卫星系统中存在一项独特现象：纳玛卡的轨道受到希伊阿卡的强烈扭矩作用，使得纳玛卡–海卫密的凌掩观测视角在后续数年内得以持续$^{[71][75][77]}$。2009 年 6 月 19 日，在巴西的皮科·多斯·迪亚斯天文台（Pico dos Dias Observatory）成功观测到一次掩食事件$^{[79]}$。
\begin{table}[ht]
\centering
\caption{Haumean系统}\label{tab_Haumea1}
\begin{tabular}{|c|c|c|c|c|}
\hline
\textbf{姓名}&\textbf{直径 (km)$^{[80]}$} & \textbf{半长轴 (km)$^{[81]}$} & \textbf{质量 (kg)$^{[81]}$} & \textbf{发现日期$^{[82]}$} \\
\hline
Haumea &  $2 322 \times 1,704 \times 1,026 $ &  & $ (4.006 \pm 0.040) \times 10^{21}$ & 2003年3月7日$^{[82]}$ \\
\hline
Hiʻiaka & $\approx310$ & 49 880 & $(1.79 \pm 0.11) \times 10^{19}$ & 2005年1月26日 \\
\hline
Namaka& $\approx170$ & 25 657 & $(1.79 \pm 1.48) \times 10^{18}$& 2005 \\
\hline
\end{tabular}
\end{table}
\subsection{碰撞族群}
海卫密（Haumea）是其碰撞族群（collisional family）中最大的成员。该族群由一组在物理性质和轨道特征上相似的天体组成，被认为是在某个更大的母体因一次撞击而破碎后形成的 $^{[37]}$。这是在跨海王星天体（TNOs）中首次被识别出的一个碰撞族群，其中包括——除海卫密及其卫星外——(55636) 2002 TX300（≈364 km）、(24835) 1995 SM55（≈174 km）、(19308) 1996 TO66（≈200 km）、(120178) 2003 OP32（≈230 km）和 (145453) 2005 RR43（≈252 km）$^{[6]}$。Brown 及其同事提出，该族群是移除海卫密冰质地幔的那次撞击的直接产物 $^{[37]}$，但另一种观点认为起源可能更加复杂：在最初的撞击中喷射出的物质首先聚合成一个海卫密的大型卫星，随后该卫星在第二次撞击中被击碎，并向外抛射其碎片 $^{[85]}$。第二种情形产生的碎片速度分布更接近该族群成员观测到的速度离散情况 $^{[85]}$。

这一碰撞族群的存在可能意味着海卫密及其“子体”可能起源于散布盘（scattered disc）。在如今稀疏的柯伊伯带中，在太阳系年龄范围内发生这种撞击的概率不足 0.1\% $^{[86]}$。而该族群不可能形成于更致密的原始柯伊伯带，因为这样紧密的天体系统会在海王星向外迁移进入柯伊伯带的过程中被破坏——这一迁移被认为是导致现今柯伊伯带低密度的原因 $^{[86]}$。因此，更可能的情形是：动态活跃的散布盘区域——其内部发生此类撞击的概率要高得多——是产生海卫密及其相关天体的起源地 $^{[86]}$。

由于该族群需要至少十亿年的时间才能扩散至当前分布范围，据信形成海卫密族群的那次撞击发生在至少十亿年前 $^{[6]}$。
\subsection{探索}
“2007年10月、2017年1月以及2020年5月，“新视野”号（New Horizons）探测器从分别为49 AU、59 AU 和 63 AU 的距离对 Haumea 进行了远距离观测$^{[19]}$。探测器的出境轨道使得能够在从地球无法获得的高相位角条件下观测 Haumea，从而得以确定 Haumea 表面的光散射特性与相位曲线行为$^{[19]}$。

一项飞掠任务如果分别在 2026 年 11 月 1 日、2037 年 9 月 23 日以及 2038 年 10 月 29 日发射，则可在 16.45 年内抵达 Haumea$^{[87]}$。Haumea 有望成为未来探测任务的目标$^{[88]}$，例如已有一项关于前往 Haumea 及其卫星（距离在 35–51 AU）探测器的初步研究$^{[89]}$。探测器的质量、动力来源以及推进系统是此类任务的关键技术领域$^{[88]}$。”
\subsection{参见}
\begin{itemize}
\item 20000 Varuna —— 一颗大型、快速自转的海王星外天体，具有拉长的椭球体形状
\item 208996 Achlys —— 另一颗大型、快速自转的海王星外天体，其拉长的椭球体形状与 Haumea 相似
\item 天文学命名惯例（Astronomical naming conventions）
\item 《我如何“杀死”冥王星，以及它该有此下场》 —— Mike Brown 于 2010 年出版的回忆录
\item 按大小排列的太阳系天体列表（List of Solar System objects by size）
\end{itemize}
\subsection{注释}
\begin{enumerate}
\item 在夏威夷的英语发音中读作 how-MAY-ə，为三个音节；根据 Brown 的学生所述，读作 HAH-oo-MAY-ə，为四个音节。$^{[1][2][3]}$
\item 假设轨道为近似圆形且偏心率可以忽略，则平均轨道速度可用完成一次轨道周长运行所需的时间 $T$ 来近似表示，其中轨道半径为其半长轴 $a$：
  [
  v \approx \frac{2\pi a}{T}.~
  ]
\item 假设 Haumea 处于流体静力平衡（hydrostatic equilibrium）的最佳拟合物理模型。$^{[11]}$
\item 基于掩星观测（occultation）并假设 Haumea 的环不对其总亮度产生贡献的模型。$^{[12]}$
\item 基于掩星观测并在“Haumea 的环对其总亮度贡献 5\%”这一上限假设之下的模型。$^{[12]}$
\item Kondratyev 与 Kornoukhov（2018）给出了以赤道坐标表示的 Haumea 北极方向，其中 $\alpha$ 为赤经（right ascension），$\delta$ 为赤纬（declination）。$^{[17]:3174}$
将赤道坐标转换为黄道坐标后，对于第一组解 $(\alpha, \delta) = (282.6^\circ, -13.0^\circ)$，得到 $\lambda \approx 282.5^\circ$、$\beta \approx 9.9^\circ$；对于第二组解 $(\alpha, \delta) = (282.6^\circ, -11.8^\circ)$，得到 $\lambda \approx 282.6^\circ$、$\beta \approx 11.1^\circ$。$^{[18]}$
黄道纬度（$\beta$）表示偏离黄道面的角度，而相对于黄道面的轨道倾角（$i$）则表示偏离黄道北极（$\beta = +90^\circ$）的角度；相对于黄道平面的倾角 $i$ 为 $\beta$ 的余角，其关系为 $i = 90^\circ - \beta$。因此，对应上述 $\beta$ 的两个数值，Haumea 的自转轴倾角相对于黄道平面分别约为 $81.2^\circ$ 与 $78.9^\circ$。
\end{enumerate}
\subsection{参考资料}
\begin{enumerate}
\item 以夏威夷女神命名的新矮行星 已存档2015-12-08在回程机(HeraldNet，2008年9月19日)
\item "DPS08网络流"。存档自原来的在2009-01-06。已检索2009-02-14。
\item 《天文学365天》。2009年3月31日。存档自原来的on 2012-02-20。已检索2009-04-11。
\item “MPEC 2010-H75: 遥远的小行星 (2010年5月14日。0 TT)”。小行星中心。2010-04-10。已存档从原来的2014-07-16。已检索2010-07-02。
\item 马克·W.Buie(2008-06-25)。“136108的轨道配合和天体测量记录”。西南研究院 (空间科学部)。已存档从原来的2011-05-18。已检索2008-10-02。
\item Ragozzine, D.;布朗，M.E.(2007)。“柯伊伯带天体2003 EL家族的候选成员和年龄估计”61”。天文杂志。134(6):2160-2167。arXiv:0709.0328。Bibcode:2007AJ....134.2160R。doi:10.1086/522334。S2CID8387493。  
\item 例如。Giovanni Vulpetti (2013)快速太阳能航行,第333页。
\item "(136108) Haumea = 2003 EL61"。小行星中心。国际天文学联合会。已存档从原件到2021年7月24日。已检索3月14日2021。
\item “喷气推进实验室小体数据库浏览器: 136108 Haumea (2003 EL61)"(2019-08-26最后一次obs)。NASA喷气推进实验室已存档从原件上2020-07-11。已检索2020-02-20。
\item “2133年6月1日左右在近日点的Haumea的Horizons批次”。JPL地平线(当rdot从负翻转为正时发生近日点。JPL SBDB通常 (错误地) 列出了2132年不受干扰的两体近日点日期)。喷气推进实验室。已存档从原始的2021-09-13。已检索9月13日2021。
\item 邓纳姆，E.T.;Desch, S.J.;普罗布斯特，L.(2019年4月)。“Haumea的形状，组成和内部结构”。天体物理学杂志。877(1): 11。arXiv:1904.00522。Bibcode:2019ApJ...877...41D。doi:10.3847/1538-4357/ab13b3。S2CID90262114。  
\item 奥尔蒂斯，J.L.;桑托斯-桑兹，P.；西卡迪，B.；贝内代蒂-罗西，G.；Bérard, D.;莫拉莱斯，N.; 等人 (2017)。“来自恒星掩星的矮行星Haumea的大小，形状，密度和环”(PDF)。自然。550(7675):219-223。arXiv:2006.03113。Bibcode:2017 Natur.550..219O。doi:10.1038/自然24051。高密度脂蛋白:10045/70230。PMID29022593。S2CID205260767。已存档(PDF)从原始的2020-11-07。已检索2020-08-19。   
\item “椭球表面积: 8.13712 × 10 ^ 6 km2"。wolframalpha.com。2019年12月20日。已存档从原件到2020年7月25日。已检索12月20日2019。
\item “椭球体体积: 1.98395 × 10 ^ 9 km3"。wolframalpha.com。2019年12月20日。已存档从原件到2020年7月25日。已检索12月20日2019。
\item Ragozzine, D.;布朗，M.E.(2009)。“矮行星Haumea = 2003 EL61的卫星的轨道和质量”。天文杂志。137(6):4766-4776。arXiv:0903.4213。Bibcode:2009AJ....137.4766R。doi:10.1088/0004-6256/137/6/4766。S2CID15310444。  
\item 桑托斯-桑兹，P.；Lellouch, E.;Groussin, O.;Lacerda, P.;穆勒，T.G.;奥尔蒂斯，J.L.;吻，C.；维勒纽斯，E.；斯坦斯伯里，J.；Duffard, R.;Fornasier, S.;乔达，L.；Thirouin, A.(2017年8月)。”“TNOs很酷”: 对跨海王星地区XII的调查。Haumea的热光曲线，2003 VS2和2003 AZ84与赫歇尔/PACS “。天文学与天体物理学。604(A95): 19。arXiv:1705.09117。Bibcode:2017A & A...604A..95S。doi:10.1051/0004-6361/201630354。S2CID 119489622。
\item Kondratyev, B.P.;科尔努克霍夫，V.S。(2018年8月)。“通过对恒星掩星和光度数据的观察来确定矮行星Haumea的身体”。皇家天文学会月刊。478(3):3159-3176。Bibcode:2018MNRAS.478.3159K。doi:10.1093/mnras/sty1321。
\item "坐标变换和银河消光计算器"。NASA/IPAC河外数据库。加州理工学院。已存档从2023年1月22日开始。已检索2月11日2023。赤道 → 黄道，J2000为春分和纪元。注意: 输入赤道坐标时，请以 “282.6d” 而不是 “282.6” 的格式指定单位。
\item 安妮·J·维比瑟；保罗·赫尔芬斯坦; 西蒙·B·波特；苏珊·D·贝内基；卡夫拉斯，J。J.;劳尔，Tod R.; 等人 (2022年4月)。“从新视野观测到的矮行星和大型KBO相位曲线的不同形状”。行星科学杂志。3(4): 31。Bibcode:2022PSJ.....3...95V。doi:10.3847/PSJ/ac63a6。95。
\item 查德威克A.特鲁希略;迈克尔·E.棕色; Kristina Barkume; Emily Shaller;大卫·L.拉比诺维茨(2007)。“2003年EL的表面61在近红外 ”。天体物理学杂志。655(2):1172-1178。arXiv:astro-ph/0601618。Bibcode:2007年4月。655.1172吨。doi:10.1086/509861。S2CID118938812。  
\item 斯诺德格拉斯，C.；携带，B.；大仲马，C.；Hainaut, O.(2010年2月)。“(136108) Haumea家族候选成员的特征”。天文学和天体物理学。511: a72。arXiv:0912.3171。Bibcode:2010A & A...511A..72S。doi:10.1051/0004-6361/200913031。S2CID62880843。  
\item 拉比诺维茨，D。L.;Schaefer, Bradley E.;舍费尔，玛莎; 图尔特洛特，苏珊娜W。(2008)。“2003年EL的年轻外观61碰撞家庭 “。天文杂志。136(4):1502-1509。arXiv:0804.2864。Bibcode:2008AJ....136.1502R。doi:10.1088/0004-6256/136/4/1502。S2CID117167835。  
\item "AstDys (136108) Haumea星历表"。意大利比萨大学数学系。已存档从原来的2011-06-29。已检索2009-03-19。
\item "HORIZONS Web-Interface"。NASA喷气推进实验室太阳系动力学。已存档从原来的2008-07-18。已检索2008-07-02。
\item "IAU命名第五矮行星Haumea"。IAU新闻稿。2008-09-17。存档自原来的on 2011-07-02。已检索2008-09-17。
\item 迈克尔·E·布朗。“2003年发现EL的电子踪迹61"。加州理工学院。已存档从原来的2006-09-01。已检索2006-08-16。
\item 巴勃罗·桑托斯·桑兹 (2008-09-26)。"La historia de Ataecina vs Haumea"(西班牙语)。infoastro.com。已存档从原始的2008-09-29。已检索2008-09-29。
\item 迈克尔·E.棕色。我如何杀死冥王星，为什么它会来第九章 “第十颗行星”
\item 杰夫·赫克特 (2005-09-21)。天文学家否认不当使用网络数据。Ne w Scientist.com。已存档从原来的2011-03-13。已检索2009-01-12。
\item “矮行星及其系统”。美国地质调查局行星命名法地名。已存档从原来的2011-06-29。已检索2008-09-17。
\item 雷切尔·考特兰 (2008-09-19)。有争议的矮行星最终命名为 “haumea” "。NewScientistSpace。已存档从原始的2008-09-19。已检索2008-09-19。
\item "圣诞老人等"。NASA天体生物学杂志。2005-09-10。从原始存档于2006-04-26。已检索2008-10-16。
\item “天文物体的命名: 小行星”。国际天文学联合会。已存档从原始的2008-12-16。已检索2008-11-17。
\item 迈克·布朗 (2008-09-17)。"矮行星: Haumea"。加州理工学院。已存档从原始的2008-09-15。已检索2008-09-18。
\item 罗伯特D.克雷格 (2004)。波利尼西亚神话手册。ABC-克里奥。第128页。ISBN 978-1-57607-894-5。已存档从原来的2023-02-08。已检索2020-11-11。
\item “新闻稿-IAU0807: IAU命名了第五个矮行星Haumea”。国际天文学联合会。2008-09-17。已存档从原来的2009-07-08。已检索2008-09-18。
\item 布朗，M.E.;Barkume，K。M、；Ragozzine, D.;Schaller, L.(2007)。“柯伊伯带中冰冷物体的碰撞族”(PDF)。自然。446(7133):294-296。Bibcode:2007 Natur.446..294B。doi:10.1038/nature05619。PMID17361177。S2CID4430027。已存档(PDF)从2020-05-04的原件。已检索2019-07-14。   
\item "建议的新字符: 管道"。已存档从原来的2022-01-29。已检索2022-01-29。
\item 柯克·米勒 (2021年10月26日)。"Unicode请求矮行星符号"(PDF)。unicode.org。已存档(PDF)从原始的2022年3月23日。已检索8月6日2022年。
\item JPL/NASA (2015年4月22日)。“什么是矮行星？”。喷气推进实验室。已存档从原始的2021-01-19。已检索2021-09-24。
\item 德博拉·安德森 (2022年5月4日)。走出这个世界: 新的天文学符号被批准为Unicode标准。unicode.org。Unicode联盟。已存档从2022年8月6日开始。已检索8月6日2022年。
\item 迈克尔·E.布朗。“最大的柯伊伯带物体”(PDF)。加州理工学院。已存档(PDF)从原始的2008-10-01。已检索2008-09-19。
\item 内斯沃尼，D; 罗伊格，F。(2001)。“跨内尼亚地区的平均运动共振第二部分: 1 : 2、3: 4和较弱的共振”。伊卡洛斯。150(1):104-123。Bibcode:2001Icar .. 150 .. 104N。doi:10.1006/icar.2000.6568。S2CID15167447。 
\item 马克·J·库什纳；迈克尔·E·布朗；霍尔曼，马修 (2002)。“经典柯伊伯带天体的长期动力学和轨道倾斜”。天文杂志。124(2):1221-1230.arXiv:astro-ph/0206260。Bibcode:2002AJ....124.1221K。doi:10.1086/341643。S2CID12641453。  
\item 拉比诺维茨，D。L.;克里斯蒂娜·巴库姆; 迈克尔·E·布朗；Roe，Henry; Schwartz，Michael; Tourtellotte，Suzanne; Trujillo，Chad (2006)。“光度学观测限制了2003年EL的大小，形状和反照率61,柯伊伯带中快速旋转的冥王星大小的物体 “。天体物理学杂志。639(2):1238-1251。arXiv:astro-ph/0509401。Bibcode:2006ApJ...639.1238R。doi:10.1086/499575。S2CID11484750。  
\item C.A.特鲁希略 & M。E.布朗 (2003年6月)。“加州理工学院广域天空调查”。地球、月球和行星。112(1-4):92-99。Bibcode:2003EM & P...92...99T。doi:10.1023/B: 月0000031929.19729.a1。S2CID189905639。  
\item 布朗，M.E.;特鲁希略，C.；拉比诺维茨，D。L.(2004)。“发现候选内奥尔特云小行星”。天体物理学杂志。617(1):645-649。arXiv:astro-ph/0404456。Bibcode:2004年4月。617 .. 645B。doi:10.1086/422095。S2CID7738201。  
\item Schwamb, M.E.;布朗，M.E.;拉比诺维茨，D。L.(2008)。“对塞德纳地区遥远人口的限制”。美国天文学会，DPS会议 #40，#38.07。40: 465。Bibcode:2008DPS....40.3807S。
\item “136108的轨道和天体测量”。www.boulder.swri.edu。已存档从原件上2020-07-13。已检索2020-07-14。
\item 法新社 (2009-09-16)。“天文学家锁定了菱形的Haumea”。欧洲行星科学大会在波茨坦举行。新闻有限公司。存档自原来的on 2009-09-23。已检索2009-09-16。
\item 斯坦斯伯里，J.；格伦迪，W.；布朗，M.；克鲁克申克，D.；斯宾塞，J.；特里林，D.；Margot, J-L.(2008)。“柯伊伯带和半人马座天体的物理特性: 斯皮策太空望远镜的约束”。超越海王星的太阳系。亚利桑那大学出版社: 161.arXiv:astro-ph/0702538。Bibcode:2008ssbn.book .. 161S。
\item 亚历山德拉C.洛克伍德; 迈克尔E.布朗; 约翰·斯坦伯里 (2014)。“长方形矮行星Haumea的大小和形状”。地球、月球和行星。111(3-4):127-137。arXiv:1402.4456v1。Bibcode:2014EM & P .. 111 .. 127L。doi:10.1007/s11038-014-9430-1。S2CID18646829。  
\item 布朗，M.E.;Bouchez, A.H、; 拉比诺维茨，D.；Sari, R.;特鲁希略，C。A.;范·达姆，M.；坎贝尔，R.；钦，J.；哈特曼，S.；约翰逊，E.；Lafon, R.;Le Mignant, D.;Stomski, P.;萨默斯，D.；Wizinowich，P。(2005)。“凯克天文台激光导星自适应光学发现和表征一颗卫星到大型柯伊伯带天体2003 EL61"(PDF)。天体物理学杂志快报。632(1):L45-L48。Bibcode:2005ApJ...632L..45B。doi:10.1086/497641。S2CID119408563。已存档(PDF)从原来的2017-12-02。已检索2018-11-04。  
\item Lacerda, P.;Jewitt, D.C.(2007)。“从它们的旋转光曲线看太阳系物体的密度”。天文杂志。133(4):1393-1408。arXiv:astro-ph/0612237。Bibcode:2007AJ....133.1393L。doi:10.1086/511772。S2CID17735600。  
\item Lellouch, E.;吻，C.；桑托斯-桑兹，P.；穆勒，T.G.;Fornasier, S.;Groussin，O.; 等人 (2010年)。“ “TNOs很酷”: 对跨海王星地区的调查II。(136108) Haumea的热光曲线。天文学和天体物理学。518: l147。arXiv:1006.0095。Bibcode:2010A & A...518L.147L。doi:10.1051/0004-6361/201014648。S2CID 119223894。
\item Fornasier, S.;Lellouch, E.;Müller, T.;桑托斯-桑兹，P.；Panuzzo, P.;吻，C.；林，T.；莫默特，M.；Bockelée-Morvan, D.；维勒纽斯，E.；斯坦斯伯里，J.；Tozzi, G.P.;Mottola, S.;德尔桑蒂，A.；克罗维西耶，J.；Duffard, R.;亨利，F.；Lacerda, P.;巴鲁奇，A.；Gicquel, A.(2013)。" “TNOs很酷”: 对跨海王星地区的调查VIII。结合Herschel PACS和SPIRE在70-500μm处对9个明亮目标的观测 (PDF)。天文学和天体物理学。555: a15。arXiv:1305.0449。Bibcode:2013A & A...555A..15F。doi:10.1051/0004-6361/201321329。S2CID 119261700。已存档 (PDF)从原来的2014-12-05。
\item “Charon: 终极冷冻中的制冰机”(新闻稿)。双子座天文台。2007年7月17日。已存档从2011年6月7日的原件。已检索2007-07-18。
\item 布朗，M.E.;夏勒，E.L.;罗伊，h.g.；拉比诺维茨，D。L.;特鲁希略，C。A.(2006)。“用哈勃太空望远镜直接测量2003 UB313的大小”(PDF)。天体物理学杂志快报。643(2):L61-L63。arXiv:astro-ph/0604245。Bibcode:2006ApJ...643L..61B。doi:10.1086/504843。S2CID16487075。已存档(PDF)从原始的2008-09-10。  
\item 皮尼拉-阿隆索，N.；布鲁内托，R.；利桑德罗，J.；吉尔-赫顿，R.；Roush, T.L.;Strazzulla, G.(2009)。“研究2003 EL61的表面，这是跨尼普顿带中最大的贫碳物体”。天文学和天体物理学。496(2):547-556。arXiv:0803.1080。Bibcode:2009A & A...496 .. 547P。doi:10.1051/0004-6361/200809733。S2CID15139257。  
\item 泰格勒，S.C.;格伦迪，W。M、；罗曼辛，W.；Consolmagno, G.J.;莫格伦，K.；维拉斯，F.(2007)。“大型柯伊伯带物体的光谱学136472 (2005 FY9) 和136108 (2003年EL61) “。天文杂志。133(2):526-530。arXiv:astro-ph/0611135。Bibcode:2007AJ....133..526T。doi:10.1086/510134。S2CID10673951。  
\item P。拉塞达; D.Jewitt & N.Peixinho (2008)。“Extreme KBO 2003 EL61的高精度光度法”。天文杂志。135(5):1749-1756.arXiv:0801.4124。Bibcode:2008AJ....135.1749L。doi:10.1088/0004-6256/135/5/1749。S2CID115712870。  
\item P。Lacerda (2009)。“极端柯伊伯带物体Haumea的时间分辨近红外光度法”。天文杂志。137(2):3404-3413。arXiv:0811.3732。Bibcode:2009AJ....137.3404L。doi:10.1088/0004-6256/137/2/3404。S2CID15210854。  
\item “奇怪的矮行星有红色斑点”。Space.com。2009年9月15日。已存档从2009年11月21日的原件。已检索2009-11-12。
\item 惊喜!矮行星Haumea有一个戒指 已存档2017-10-22在回程机天空和望远镜，2017年10月13日。
\item Kondratyev, B.P.;科尔努克霍夫，V.S。(2020年10月)。“旋转三轴引力体周围的环的长期演化”。天文学报告。64(10):870-875。Bibcode:2020年阿雷普。64。870k。doi:10.1134/S1063772920100030。
\item 冬天，O.C.;Borderes-Motta, G.;里贝罗，T。(2019年)。“关于矮行星Haumea周围的环的位置”。皇家天文学会月刊。484(3):3765-3771。arXiv:1902.03363。doi:10.1093/mnras/stz246。S2CID119260748。 
\item K。Chang (2007年3月20日)。“将旧碰撞的线索拼凑在一起，一个冰球”。纽约时报。已存档从2014年11月12日的原件。已检索2008-10-12。
\item 布朗，M.E.；范·达姆，M。A.;Bouchez, A.H.; Le Mignant, D.;坎贝尔，R。D.;Chin, J.C.Y.;康拉德，A.；哈特曼，S.K.;约翰逊，E。M、；Lafon, R.E.;拉比诺维茨，D。L.拉比诺维茨; Stomski，P。J.小; 萨默斯，D。M、；特鲁希略，C。A.;Wizinowich，P。L.(2006)。“最大的柯伊伯带物体的卫星”(PDF)。天体物理学杂志。639(1):L43-L46。arXiv:astro-ph/0510029。Bibcode:2006ApJ...639L..43B。doi:10.1086/501524。S2CID2578831。已存档(PDF)从原来的2013-11-03。已检索2011-10-19。  
\item K。M。Barkume; M.E.布朗 & E。L.Schaller (2006)。“柯伊伯带天体2003 EL卫星上的水冰61”。天体物理学杂志快报。640(1):L87-L89。arXiv:astro-ph/0601534。Bibcode:2006ApJ...640L..87B。doi:10.1086/503159。S2CID17831967。  
\item 格林，丹尼尔·W.E.(2005年12月1日)。"Iauc 8636"。已存档从2018年3月12日的原件开始。
\item Ragozzine, D.;布朗，M.E.;特鲁希略，C。A.;夏勒，E.L.(2008)。2003 EL61卫星系统的轨道和质量。2008年AAS DPS会议。美国天文学会公报。第40卷.第462页。Bibcode:2008DPS....40.3607R。
\item Ragozzine, D.;布朗，M.E.(2009)。“矮行星Haumea = 2003 EL的卫星的轨道和质量61”。天文杂志。137(6):4766-4776。arXiv:0903.4213。Bibcode:2009AJ....137.4766R。doi:10.1088/0004-6256/137/6/4766。S2CID15310444。  
\item 大仲马，C.；携带，B.；赫斯特罗弗，D.；梅林，F.(2011)。“(136108) Haumea的高对比度观测”。天文学与天体物理学。528: a105。arXiv:1101.2102。Bibcode:2011A & A...528A.105D。doi:10.1051/0004-6361/201015011。已检索8月24日2024。
\item "IAU通告8949"。国际天文学联合会。2008年9月17日。存档自原来的2009年1月11日。已检索2008-12-06。
\item “Haumea和Namaka的共同事件”。已存档从原来的2009-02-24。已检索2009-02-18。
\item L.-A。A.麦克法登; P。R。魏斯曼; T.V。约翰逊 (2007)。太阳系百科全书。学术出版社。ISBN  978-0-12-088589-3。
\item Fabrycky, D.C.;霍尔曼，M。J.;Ragozzine, D.;布朗，M.E.;李斯特，T。A.;Terndrup, D.M、；Djordjevic, J.;Young, E.F.;年轻，我。A.;豪厄尔，R。R。(2008)。2003年EL的共同事件61和它的内部卫星。2008年AAS DPS会议。美国天文学会公报。第40卷.第462页。Bibcode:2008DPS....40.3608F。
\item M。布朗 (2008年5月18日)。"月影星期一 (已修复)"。迈克·布朗的行星。已存档从2008年10月1日起。已检索2008-09-27。
\item 博尔托莱托，A.；Saito, R.K。(2010)。“观察来自巴西的跨海王星物体Haumea和Namaka的相互事件”。国际天文学联合会会议记录。269:189-192.Bibcode:2010IAUS..269..189B。doi:10.1017/S1743921310007404。已检索8月24日2024。
\item 奥尔蒂斯，J.L.;桑托斯-桑兹，P.；西卡迪，B.；贝内代蒂-罗西，G.；Bérard, D.;莫拉莱斯，N.；Duffard, R.;布拉加-里巴斯，F.；霍普，美国；里斯，C.；Nascimbeni, V.(2017年10月)。“来自恒星掩星的矮行星Haumea的大小，形状，密度和环”。自然。550(7675):219-223。arXiv:2006.03113。Bibcode:2017 Natur.550..219O。doi:10.1038/自然24051。ISSN0028-0836。PMID29022593。S2CID205260767。已存档从原始的2022-06-23。已检索2021-07-08。   
\item Ragozzine, D.;布朗，M.E.(2009-06-01)。“矮行星Haumea (2003 El61) 的卫星的轨道和质量”。天文杂志。137(6):4766-4776。arXiv:0903.4213。Bibcode:2009AJ....137.4766R。doi:10.1088/0004-6256/137/6/4766。ISSN0004-6256。S2CID15310444。已存档从原始的2021-05-09。已检索2021-07-08。  
\item "深度 | Haumea"。NASA太阳系探索。2017年11月14日.已存档从原始的2021年6月29日。已检索7月8日，2021。
\item Proudfoot，本杰明; Ragozzine，达林 (2019年5月)。“矮行星Haumea家族的形成建模”。天文杂志。157(6): 230。arXiv:1904.00038。Bibcode:2019AJ....157..230P。doi:10.3847/1538-3881/ab19c4。S2CID90262136。 
\item Proudfoot，本杰明; Fernández-valenzuela，Estela; Stansberry，John; Kiss，Csaba; Ragozzine，达林; 弗雷泽，韦斯利; 派克，迷迭香; 皮尼拉-阿隆索，诺埃米 (2024年12月)。“候选Haumea家庭成员的近红外调查”。天文杂志。168(6): 269。Bibcode:2024AJ....168..269P。doi:10.3847/1538-3881/ad864e。S2CID274155364。 
\item 施利希廷，H. E.；Sari, R.(2009)。“Haumea的碰撞家庭的创造”。天体物理学杂志。700(2):1242-1246。arXiv:0906.3893。Bibcode:2009年4月。700.1242s。doi:10.1088/0004-637X/700/2/1242。S2CID19022987。  
\item 利维森，h.f.；莫比德利，A.；Vokrouhlický, D.;Bottke, W.F.(2008)。“在2003年EL的分散圆盘原点上61碰撞族-碰撞在小物体动力学中的重要性的一个例子 “。天文杂志。136(3):1079-1088。arXiv:0809.0553。Bibcode:2008AJ....136.1079L。doi:10.1088/0004-6256/136/3/1079。S2CID10861444。  
\item 麦格拉纳汉，R.；萨根，B.；鸽子，G.；图洛斯，A.；Lyne, J.E.;埃默里，J。P。(2011)。“对跨尼普顿天体的任务机会的调查”。英国星际学会杂志。64:296-303。Bibcode:2011JBIS...64..296M。
\item 乔尔·庞西; 乔迪·丰德卡巴·贝加; 弗雷德·费雷辛布; 文森特·马蒂诺塔 (2011年)。“对haudmean系统中轨道飞行器的初步评估: 行星轨道飞行器能以多快的速度到达如此遥远的目标？”。宇航学报。68(5-6):622-628.Bibcode:2011年AcAau。。68 .. 622P。doi:10.1016/j.actaastro.2010.04.011。
\item 保罗·吉尔斯特:到Haumea的快速轨道器 已存档2015-09-23在回程机。半人马座梦想-Tau Zero基金会的新闻。2009年7月14日，2011年1月15日检索
\end{enumerate}
\subsection{外部链接}
\begin{itemize}
\item (136108) Haumea、hi'iaka和Namaka在约翰斯顿Archive.com (2014年9月21日更新)
\item 2009国际天文年 播客: 矮行星Haumea (Darin Ragozzine)
\item 2011年6月10日看到的Haumea由Mike Brown使用4.20 m (165 in)WHT/~ 0:30-3:30浸入Haumea + Namaka的亮度当Namaka穿过Haumea时(外月hi'iaka混合在图像中，但每4.5小时旋转一次，并增加了一点变化)
\item 在接下来的350万年中，Haumea与海王星间歇性7:12共振的动画
\end{itemize}